\subsection{\texorpdfstring{Integrable bundles in characteristic \(p \neq 0\)}{Integrable bundles in characteristic p nonzero}}

In this number, we assume K has characteristic \(p \neq 0\) . We denote by X a K-analytic manifold locally of finite dimension.

9.4.1 ("p-th powers"). Let \(\xi\) be a vector field on an open subset U of X. There exists a vector field \(\xi^p\) and only one on U such that we have

\[
\mathrm{D} _ {\xi^ {p}} (f) = (\mathrm{D} _ {\xi}) ^ {p} (f) = \underbrace {\mathrm{D} _ {\xi} (\mathrm{D} _ {\xi} (\dots (\mathrm{D} _ {\xi} (f)) . . .))} _ {p \text { times }}\tag{1}
\]

for every analytic function f defined in an open subset of U.

If \(\varphi\) is an analytic function on U, we have

\[
(\varphi \xi) ^ {p} = \varphi^ {p} \xi^ {p} + (\mathrm{D} _ {\varphi \xi}) ^ {p - 1} (\varphi). \xi\tag{2}
\]

and if \(\eta\) is a vector field on U, we have

\[
[ \xi^ {p}, \eta ] = \operatorname{ad}(\xi) ^ {p} (\eta) = [ \xi , [ \xi , \dots , [ \xi , \eta ] \dots ] ].\tag{3}
\]

9.4.2 ("Jacobson's identity"). Let L be the free \(F_{p}\)-Lie algebra (LIE, II, § 2, no. 2) on a set \(\{x, y\}\) with two elements. There exists one and only one element \(\Lambda_{p}(x, y)\) of L such that we have

\[
(x + y) ^ {p} = x ^ {p} + y ^ {p} + \Lambda_ {p} (x, y)\tag{4}
\]

in the enveloping algebra of L. Examples:

\[
\Lambda_ {2} (x, y) = [ x, y ]; \quad \Lambda_ {3} (x, y) = [ x, [ x, y ] ] - [ y, [ x, y ] ].
\]

If \(\xi\) and \(\eta\) are vector fields on X, we have

\[
(\xi + \eta) ^ {p} = \xi^ {p} + \eta^ {p} + \Lambda_ {p} (\xi , \eta).\tag{5}
\]

9.4.3. Example. --- Take \(X = K\) (resp. \(K^*\)) and denote by x the canonical map \(X \to K\). Let \(\xi\) (resp. \(\eta\)) be the vector field \(\partial/\partial x\) (resp. \(x\cdot\partial/\partial x\)); it is invariant under translations of the additive (resp. multiplicative) group X. We have

\[
\xi^ {p} = 0 \quad \text { and } \quad \eta^ {p} = \eta .\tag{6}
\]

9.4.4 (“Integrable bundles”). Let F be a vector subbundle of T(X). We say that F is integrable if, for every \(x \in X\) , there exists a coordinate system (\(\zeta^{1}, \ldots, \zeta^{n}\)) of X at x and an integer \(m \leqslant n\) such that \((\partial/\partial\zeta^{1}, \ldots, \partial/\partial\zeta^{m})\) is a frame of F in a neighborhood of x.

For F to be integrable, it is necessary and sufficient that it satisfy the following condition:

For every open set U of X, \(\mathcal{S}_{\mathbb{F}}^{\omega}(\mathbf{U})\) is a Lie subalgebra of \(\mathcal{S}_{\mathrm{T}(\mathbf{X})}^{\omega}(\mathbf{U})\) stable under the operation \(\xi \mapsto \xi^{p}\) .

